\section{Ethical Considerations}
\label{sec:ethics}

% TRACE: S9-C01 | configs/data_governance.yaml | dataset_construction, access_policy, public_release_status; answer/data_provenance/source_registry.json | $.sources; answer/THIRD_PARTY_NOTICES.md | Dataset source
ViPragSent is handled as private, non-commercial research data.  Its fixed gold corpus contains retained local ViSoBERT-export records and author-created, context-augmented derivatives based on VIVID idiom/proverb materials.  Raw and processed source text are not redistributed, and source authorization and license documentation are pending project action.  The manuscript reports aggregate results, task definitions, and provenance records without publishing raw text or source-text examples.  External evaluation datasets are not components of ViPragSent; they are used only for the reported diagnostic and are not redistributed by this project.

% TRACE: S9-C01a | answer/results/p1_source_stratified_sensitivity.json; configs/data_governance.yaml | P1 result scope and pending VIVID governance | results do not imply access or redistribution rights
The source-stratified P1 analysis does not change that governance boundary. The existence of aggregate results for a VIVID-labelled stratum is not evidence of VIVID license clearance, research-use authorization, unrestricted access, or a right to redistribute source or derivative text.

% TRACE: S9-C02 | configs/data_governance.yaml | access_policy, public_release_status; answer/LICENSE | Sections 2--5
This boundary is intended to reduce privacy and contextual-integrity risks associated with sensitive user-generated language.  Even when text is accessible to researchers, reproducing a post, attributing it to an individual, or exposing a rare linguistic pattern can increase identifiability or cause contextual harm.  The project accordingly limits the paper to aggregate reporting and preserves data-access decisions outside the anonymous manuscript.  Any future controlled-access or release mechanism should be assessed for governance and re-identification risk before it is implemented.

% TRACE: S9-C03 | answer/results/annotation_agreement.json | $.reviewers, $.protocol_note; docs/annotation_guidelines_v1.md | Annotation workflow; answer/reproducibility/verification_manifest.json | integrity_scope
The annotation protocol uses two reviewers and a fixed adjudicator whose post-adjudication labels define the recorded gold splits.  This process does not eliminate ambiguity in pragmatic interpretation, nor does it ensure that all social groups, dialects, or communicative norms are represented equally.  Misclassification of irony, sarcasm, implicit sentiment, or other pragmatic behavior could misrepresent speakers or amplify unequal treatment when used in moderation, profiling, customer-service, or surveillance settings.  The reported systems should therefore not be used as sole decision makers about individuals.  The artifact-level traceability provided here supports inspection of recorded evidence, but it does not remove these representational risks or replace governance review in a deployment setting.
